\section{Related Work}
\label{sec:related-work}

% Facts on study designs were extracted from the full texts (arXiv HTML/PDF,
% 2026-10-25). See notes/literature-review.md for the search protocol.

\subsection{Energy Efficiency of LLM-Generated Code}

\paragraph{Single-language studies.}
Most empirical work on the energy efficiency of LLM-generated code targets Python on a
single platform. Islam et al.~\cite{islam2025evaluating} compare code generated by 20 LLMs
for 878 LeetCode problems with canonical human solutions on an Intel Xeon server, using
RAPL via \texttt{perf}, and find LLM solutions to be less energy-efficient on average, with
large differences between models and algorithmic categories.
Cappendijk et al.~\cite{cappendijk2025prompting} and Ashraf et al.~\cite{ashraf2025slm}
study the effect of prompt modifications and of model size, respectively; both report that
outcomes depend on the specific combination of model and problem.
Vartziotis et al.~\cite{vartziotis2024learn} propose a ``green capacity'' metric for
comparing code assistants.

\paragraph{Hardware platforms.}
Cursaru et al.~\cite{cursaru2024codellama} measure code generated by a single model
(Code Llama) for C++, JavaScript, and Python on a Raspberry~Pi with an external power
monitor and observe that efficiency depends strongly on language and problem.
Apsan et al.~\cite{apsan2026whereare} are, to our knowledge, the only study to execute
LLM-generated code on three heterogeneous platforms (an Intel Xeon server, an Intel Core~i9
PC, and an ARM Cortex-A72 Raspberry~Pi), comparing six LLMs and four prompting techniques
against human and expert-written Python code. Each platform is analysed separately, using
aligned rank transform (ART) ANOVA with LLM and solution type as factors. The authors
observe that the most efficient LLM differs between platforms, but this observation is not
tested statistically, and the study is restricted to Python and nine problems.

\paragraph{Multiple languages.}
Solovyeva et al.~\cite{solovyeva2025powerhungry} evaluate three LLMs (GitHub Copilot,
GPT-4o, o1-mini) on 53 hard LeetCode problems in Python, Java, and C++ on an Apple~M3
laptop and an Intel PC, measuring CPU energy with \texttt{powermetrics} and RAPL. They use
pairwise Mann--Whitney~U tests and Spearman correlations between platforms; C++ programs
are compiled without optimisation (\texttt{-O0}). Interactions between model, language, and
platform are not modelled.
Rekhi et al.~\cite{rekhi2026prompting} evaluate eight prompting strategies across ten LLMs
for Python and C++ on a single Intel cloud platform and note qualitatively that the effect
of prompting is inconsistent across models and languages, without inferential statistics.

\paragraph{Secondary studies.}
A recent systematic literature review of 18 primary studies~\cite{masthanali2026slr}
concludes that evidence is ``limited and highly heterogeneous'', that most studies evaluate
generated code on a single hardware configuration, that few use direct hardware-based
measurement, and explicitly calls for cross-platform evaluations.

\subsection{Runtime Efficiency Benchmarks}
A parallel line of work proposes benchmarks for the runtime efficiency (execution time and
memory, not energy) of generated code, e.g., EffiBench~\cite{huang2024effibench},
EvalPerf~\cite{liu2024evalperf}, ENAMEL~\cite{qiu2025enamel}, and
COFFE~\cite{peng2025coffe}. EffiBench-X~\cite{qing2025effibenchx} extends this to six
languages, including Python, Java, and C++, and reports language-specific differences in
the relative efficiency of LLMs. These benchmarks are candidate task sources for our study
(\cref{sec:setup}), but they measure on a single platform and do not report energy.

\subsection{Energy Efficiency of Programming Languages}
Pereira et al.~\cite{pereira2017energy,pereira2021ranking} rank up to 27 languages by
energy, time, and memory on the Computer Language Benchmarks Game and show that the
fastest language is not always the most energy-efficient. Their results concern
human-written, highly optimised programs on a single x86 machine; whether similar
relations hold for LLM-generated code and across architectures is open.

\subsection{Research Gap}
\Cref{tab:related-work} summarises the closest studies. Three observations motivate our
work:
\begin{enumerate}[leftmargin=*]
  \item \emph{No study crosses multiple languages with architecturally heterogeneous
        platforms.} Multi-language studies use one or two x86/Apple
        platforms~\cite{solovyeva2025powerhungry,rekhi2026prompting}; the only
        three-platform study is Python-only~\cite{apsan2026whereare}.
  \item \emph{Interaction effects are not tested.} Several studies report that results
        ``depend on'' model, language, or
        platform~\cite{apsan2026whereare,rekhi2026prompting,cursaru2024codellama,
        cappendijk2025prompting}, but none models these dependencies as interaction
        terms or tests the stability of LLM rankings.
  \item \emph{Measurement and statistics are heterogeneous}~\cite{masthanali2026slr}:
        several studies report averages without inferential statistics, and C++ is
        sometimes compiled without optimisation.
\end{enumerate}
Our study addresses these points with a full factorial design (LLM $\times$ language
$\times$ platform), a mixed-effects analysis with explicit interaction terms, and a test of
rank stability across languages and platforms.

\begin{table}[t]
  \centering
  \caption{Closest empirical studies on the energy efficiency of LLM-generated code.
    Ext.\ = external power meter; RAPL = software-based (Intel RAPL).}
  \label{tab:related-work}
  \small
  \setlength{\tabcolsep}{4pt}
  \begin{tabular}{@{}lclllc@{}}
    \toprule
    Study & LLMs & Languages & Platforms & Measurement & Interact. \\
    \midrule
    Cursaru et al.~\cite{cursaru2024codellama}        & 1  & C++, JS, Py    & RPi                   & Ext.            & \TBD \\
    Cappendijk et al.~\cite{cappendijk2025prompting}  & 5  & Py             & 1 (x86)               & RAPL            & no \\
    Islam et al.~\cite{islam2025evaluating}           & 20 & Py             & 1 (x86 server)        & RAPL            & no \\
    Solovyeva et al.~\cite{solovyeva2025powerhungry}  & 3  & Py, Java, C++  & Mac (M3), PC (x86)    & powermetrics, RAPL & no \\
    Apsan et al.~\cite{apsan2026whereare}             & 6  & Py             & Server, PC, RPi       & EnergiBridge, Ext. & partly$^a$ \\
    Rekhi et al.~\cite{rekhi2026prompting}            & 10 & Py, C++        & 1 (x86 cloud)         & RAPL            & no \\
    \midrule
    \textbf{This study} & \TBD & Py, Java, C++ & PC, Server, RPi & \TBD & yes \\
    \bottomrule
    \multicolumn{6}{@{}p{0.95\linewidth}@{}}{\footnotesize $^a$ LLM $\times$ prompt
      interaction tested within each platform; no interaction with platform.}
  \end{tabular}
\end{table}

\TODO{Verify Cappendijk et al.\ platform/tool and Cursaru et al.\ statistics from full
texts. Re-run the arXiv/Scopus search shortly before submission.}
